\chapter*{Abstract}
\addcontentsline{toc}{chapter}{Abstract}

PhenoHive is a portable, low-cost station that monitors a single plant continuously. It runs in the LBIR1251 plant biology course at UCLouvain, where about 20 student groups carry out month-long experiments in parallel, most of them in their own homes rather than a controlled lab. A teaching platform like this has to satisfy two demands at once: the sensor data must be accurate enough to interpret in physiological terms, and the hardware has to keep running unattended for weeks. The version I inherited fell short on three fronts. Its codebase was monolithic, so extending or maintaining it was awkward. The sensors were too basic and hadn't been validated against references, which left derived quantities such as vapor pressure deficit and spectral band assignments impossible to trust. And it coped badly with the network and power interruptions that are routine in a home.

In order to fill the gap in measurements, I calibrated the environmental and light sensors in a controlled laboratory. These corrections drop the uncertainty on measurements down to a level that allows the students to be confident in the data and will allow them to interpret the measurements in physiological terms. To achieve this, I characterized the light spectrum across all of it's channels by using a monochromator and I calibrated the environmental sensor by using a controlled chamber. The deployment also changed how the station stores and pushes data to the server. It now always writes each measurement locally first in order to not lose any data during potential network outages. Students can also connect the device to their WiFi network through a captive portal instead of a terminal, remote management runs over a private network overlay and every team's dashboard is automatically created the first time their station is connected to the server. My testing on physical hadware validated three critical things. First that a fresh station can start pushing live data to the server within five minutes. Second it that a long disconnection from the server doesn't create any data loss. And third is that a student's data is never accessible to another student from another group. 

